\chapter{Methodology}
\label{chap:Implementation}

This chapter describes the testbed built to study provider-selective distillation and
the methods evaluated on it. Section~\ref{imp:sec:overview} gives an overview of one
training run. Section~\ref{imp:sec:tasks} describes the three task domains and how
answers are checked. Section~\ref{imp:sec:providers} introduces the student, the
provider pool, the prices and the degraded provider. Section~\ref{imp:sec:update}
specifies how the student learns from a black-box provider, and
Section~\ref{imp:sec:selectors} specifies the selection policies and baselines.
Section~\ref{imp:sec:marginal} describes how marginal value is measured to test H1 and
H2. Section~\ref{imp:sec:config} lists the experimental configuration,
Section~\ref{imp:sec:verification} reports how the testbed was verified, and
Section~\ref{imp:sec:limitations} states the limitations of the design.

The notation follows Section~\ref{sec:Problem Statement}. The student is $\pi_\theta$,
the providers are $M_1, \dots, M_K$ with prices $c_1, \dots, c_K$, prompts come from
domains $\mathcal{X}_1, \dots, \mathcal{X}_D$, an arm is a pair $a = (k, d)$, each round
uses a batch of $b$ prompts, and the total budget is $B$.

\section{Overview}
\label{imp:sec:overview}

A training run alternates between choosing an arm and learning from it.
Figure~\ref{imp:fig:round} shows one round. The selector picks an arm $(k, d)$. The
student answers $b$ prompts from domain $d$, sampling $G$ responses to each. Every
response is reduced to a canonical final answer, and each response is rewarded by how
well its answer agrees with the answers that provider $k$ gave to the same prompt. The
student then takes one policy-gradient step, and the budget is charged $c_k$ for every
prompt on which the provider was consulted. Every ten rounds the student is evaluated
on a held-out selection set, and the result is passed back to the selector.

\begin{figure}[htbp]
\centering
\begin{tikzpicture}[
    box/.style={draw, rounded corners=2pt, align=center, font=\small,
                minimum height=1.1cm, text width=2.9cm},
    arr/.style={-{Stealth[length=2.2mm]}, thick}]
  \node[box] (sel)  at (0, 0)    {Selector\\chooses $(k, d)$};
  \node[box] (roll) at (4.2, 0)  {Student samples\\$G$ responses to\\$b$ prompts from $\mathcal{X}_d$};
  \node[box] (can)  at (8.4, 0)  {Canonical answers\\$\kappa_d(y)$};
  \node[box] (filt) at (8.4, -2.4) {Pre-query filter:\\skip prompts with\\identical answers};
  \node[box] (rew)  at (4.2, -2.4) {Query cached answers\\of $M_k$, charge $c_k$;\\agreement reward};
  \node[box] (upd)  at (0, -2.4)   {GRPO step on\\LoRA adapter};
  \node[box, text width=6.5cm] (ev) at (2.1, -4.6)
       {Every $L = 10$ rounds: greedy evaluation\\on the selection set, accuracies $\alpha_d$};
  \draw[arr] (sel) -- (roll);
  \draw[arr] (roll) -- (can);
  \draw[arr] (can) -- (filt);
  \draw[arr] (filt) -- (rew);
  \draw[arr] (rew) -- (upd);
  \draw[arr] (upd.south) -- ++(0, -0.55) -| (ev.north);
  \draw[arr] (ev.west) -- ++(-1.0, 0) |- (sel.west);
\end{tikzpicture}
\caption{One round of provider-selective distillation. The selector is consulted and
updated once per block of $L$ rounds.}
\label{imp:fig:round}
\end{figure}

Algorithm~\ref{imp:alg:run} gives the complete procedure. The components are described
in the sections that follow.

\begin{algorithm}[htbp]
\caption{One training run}
\label{imp:alg:run}
\begin{algorithmic}[1]
\Require student $\pi_\theta$ with a fresh LoRA adapter, providers' cached answers,
  prices $c_k$, budget $B$, round cap $T_{\max}$, block length $L$
\State $t \gets 0$, spent $\gets 0$
\State $\alpha \gets$ \Call{Evaluate}{$\pi_\theta$} \Comment{per-domain selection accuracy}
\While{$t < T_{\max}$ and some arm is affordable}
  \State $(k, d) \gets$ \Call{Select}{$\alpha$} \Comment{only arms with $b\,c_k \le B - \text{spent}$}
  \State $\alpha_{\mathrm{before}} \gets \alpha$, cost $\gets 0$
  \For{$L$ rounds, while $(k, d)$ is affordable and $t < T_{\max}$}
    \State draw the next $b$ prompts $x_1, \dots, x_b$ from $\mathcal{X}_d$
    \State sample $y_{i,1}, \dots, y_{i,G} \sim \pi_\theta(\cdot \mid x_i)$ for each $i$
    \State keep prompts whose canonical answers $\kappa_d(y_{i,g})$ are not all equal
    \State reward each kept response against $M_k$'s cached answers (Eq.~\ref{imp:eq:reward})
    \State take one GRPO step on the kept groups (Eq.~\ref{imp:eq:loss})
    \State spent $\gets$ spent $+\, c_k \cdot$ (number of kept prompts); $t \gets t + 1$
    \State save a checkpoint when progress first reaches 25\%, 50\% or 75\%
  \EndFor
  \State $\alpha \gets$ \Call{Evaluate}{$\pi_\theta$}
  \State \Call{UpdateSelector}{$(k, d)$, $\alpha_{\mathrm{before}}$, $\alpha$, cost}
\EndWhile
\State evaluate on the full selection and test sets
\end{algorithmic}
\end{algorithm}

\section{Tasks and Answer Checking}
\label{imp:sec:tasks}

\subsection{Domains}
\label{imp:sec:domains}

Three single-turn domains are used, each with answers that can be checked by a
program. Table~\ref{imp:tab:domains} summarises them. Every domain is split into a
prompt pool of 2,000 items, on which the student trains and for which provider
responses are cached, a selection set of 500 items, which supplies the selector's
feedback, and a test set of 500 items, which is used only for final results. The three
splits are disjoint within each domain.

\begin{table}[htbp]
\centering
\small
\begin{tabular}{@{}lllp{5.4cm}@{}}
\toprule
\textbf{Domain} & \textbf{Source} & \textbf{Answer} & \textbf{Filtering} \\
\midrule
Countdown & Countdown-Task-GOLD & equation & three-number puzzles only; pool puzzles that also
  appear in the test configuration removed \\
MATH & MATH \autocite{hendrycks_math_2021} & expression & difficulty levels 3 to 5, all
  seven subjects \\
Code & KodCode-V1 \autocite{xu_kodcode_2025} & function & single function, medium or
  hard, no fixtures, low similarity to public benchmarks, reference solution verified \\
\bottomrule
\end{tabular}
\caption{Task domains. Each has a 2,000-prompt pool, a 500-prompt selection set and a
500-prompt test set.}
\label{imp:tab:domains}
\end{table}

\paragraph{Countdown.} Given three integers and a target, the model must write an
arithmetic expression that uses each integer exactly once and evaluates to the target,
inside \texttt{<answer>} tags. Puzzles come from the dataset released with GOLD
\autocite{patino_gold_2025}.\footnote{\url{https://huggingface.co/datasets/HuggingFaceTB/Countdown-Task-GOLD}}
The pool is drawn from its 80,000-puzzle training configuration and the selection and
test sets from its 10,000-puzzle test configuration. Four-number puzzles were excluded
after screening: the base student solved 0.8\% of them, and only 6\% of its groups of
eight responses contained both a correct and an incorrect answer, so reinforcement
learning would have had almost nothing to learn from. On three-number puzzles the
corresponding figures were 6\% and 32\%.

\paragraph{MATH.} Competition mathematics problems at difficulty levels 3 to 5
\autocite{hendrycks_math_2021}, of which there are 5,586 in the training split and 3,669
in the test split. The pool and selection set come from the training split and the
test set from the test split. Prompts ask for step-by-step reasoning and a final answer
in \verb|\boxed{}|.

\paragraph{Code.} Python function-writing problems from KodCode-V1
\autocite{xu_kodcode_2025}, each with a pytest test file. Of its 484,097 problems, 65,087
remain after keeping instruction-style problems with a single target function, medium
or hard difficulty, no pytest fixtures or parametrisation, and a reported similarity
below 0.8 to existing benchmarks. Each problem's reference solution was then run against
its tests in the sandbox described below, and only problems whose reference passes were
kept (3,608 of 3,900 sampled candidates, 92.5\%). The prompt states the required
function signature.

\paragraph{Domains considered and rejected.} GSM8K \autocite{cobbe_gsm8k_2021} was the
original mathematics domain, but Qwen2.5-7B, Qwen2.5-14B, Qwen2.5-32B and Phi-4 all
answered between 94 and 97\% of screening prompts correctly, which leaves almost no
difference between providers for a selector to exploit. MBPP \autocite{austin_mbpp_2021} was the original
code domain, but it contains only 974 problems in total, too few for a 2,000-prompt pool.

\subsection{Canonical answers and agreement}
\label{imp:sec:canonical}

Providers are black boxes, so the student can only be compared with a provider through
the text each produces. Comparing whole responses would reward imitation of style. The
testbed instead reduces every response $y$ to a canonical final answer
$\kappa_d(y)$, or to $\varnothing$ when no valid answer can be extracted, and two
responses agree when their canonical answers are equal. The canonical form is chosen so
that different but equivalent answers compare equal.

\begin{description}[itemsep=3pt,topsep=4pt]
  \item[Countdown.] The expression in the last \texttt{<answer>} tag, with any trailing
    ``$= v$'' removed and LaTeX operators and delimiters normalised. If it uses exactly
    the given integers, $\kappa$ is its value rounded to four decimal places; otherwise
    $\kappa = \varnothing$. Two different expressions with the same value therefore
    agree.
  \item[MATH.] The content of the last \verb|\boxed{}|, parsed and normalised by the
    Math-Verify library,\footnote{\url{https://github.com/huggingface/Math-Verify}} so
    that, for example, $0.5$ and $\frac{1}{2}$ have the same canonical form.
  \item[Code.] The function's behaviour on the test inputs. The longest Python code block
    is extracted and run in the sandbox, and $\kappa$ is the tuple of values the target
    function returns, or the exception it raises, on the inputs used by the tests. The
    expected outputs are not used, so two programs agree when they behave identically,
    whether or not either is correct.
\end{description}

Correctness, which is needed only for evaluation and for the provider-free control, is
checked separately: a Countdown answer must equal the target, a MATH answer must be
judged equivalent to the reference by Math-Verify, and a program must pass every test.

\subsection{Sandbox and harness rules}
\label{imp:sec:sandbox}

Generated programs run in a bubblewrap sandbox\footnote{\url{https://github.com/containers/bubblewrap}}
with a read-only file system, no network access, a 4\,GB memory limit and time limits
of 5 seconds to load a program and 2 to 3 seconds per test. Before running, the code is
reduced to its module-level imports, definitions and assignments, which removes the
self-tests, print statements and input-reading code that models often append, and
\texttt{from typing import *} is prepended so that type hints such as
\texttt{List[int]} do not fail for lack of an import. These rules were added after
screening showed well-formed functions being scored as missing answers, and they apply
identically to the student and to every provider.

\section{Student and Providers}
\label{imp:sec:providers}

\subsection{Student}

The student is Qwen2.5-1.5B-Instruct \autocite{qwen_qwen25_2024}, trained through a
LoRA adapter \autocite{hu_lora_2022} of rank 16 on every linear layer of the attention
and feed-forward blocks. A rank of 64 was measured to add only 1.5\% to the time per
step, so cost does not constrain the choice of rank.

\subsection{Provider pool}

Nine open-weight models were screened as candidate providers on 200 selection prompts
per domain (Chapter~\ref{chap:Evaluation}), and five were kept. Two members of the
screened set were dropped because Phi-4 matched or exceeded them in every domain at an
equal or lower price (Qwen2.5-14B and Qwen2.5-32B), one because it repeated itself until
the length limit on Countdown and was therefore expensive (Granite-3.3-8B), and one
because it gave no advantage in its own speciality (Qwen2.5-Math-7B, which scored 0.73
on MATH against 0.72 for the general Qwen2.5-7B, and produced no valid Countdown
answers). Table~\ref{imp:tab:pool} lists the final pool.

\begin{table}[htbp]
\centering
\small
\setlength{\tabcolsep}{5pt}
\begin{tabular}{@{}llrrccc@{}}
\toprule
 & & & & \multicolumn{3}{c}{\textbf{Ceiling $\mu_{k,d}$}} \\
\cmidrule(l){5-7}
\textbf{Provider} & \textbf{Family} & \textbf{Params} & \textbf{Price $c_k$} &
  \textbf{Countdown} & \textbf{MATH} & \textbf{Code} \\
\midrule
Phi-4 \autocite{abdin_phi4_2024}              & Phi  & 14.7B & 2.92 & 0.874 & 0.804 & 0.686 \\
Qwen2.5-7B \autocite{qwen_qwen25_2024}        & Qwen & 7.6B  & 1.00 & 0.810 & 0.750 & 0.580 \\
Qwen2.5-Coder-7B \autocite{hui_qwen25coder_2024} & Qwen & 7.6B & 0.98 & 0.620 & 0.678 & 0.612 \\
Qwen2.5-3B \autocite{qwen_qwen25_2024}        & Qwen & 3.1B  & 0.50 & 0.538 & 0.690 & 0.462 \\
Phi-4-mini \autocite{microsoft_phi4mini_2025} & Phi  & 3.8B  & 0.84 & 0.616 & 0.618 & 0.558 \\
Phi-4, degraded                               & Phi  & 14.7B & 2.92 & 0.614 & 0.532 & 0.492 \\
\midrule
Student, before training & Qwen & 1.5B & -- & 0.195 & 0.414 & 0.438 \\
\bottomrule
\end{tabular}
\caption{The provider pool. The ceiling $\mu_{k,d}$ is the accuracy of the provider's
most common answer over its four cached responses, on the 500-prompt selection set
(Section~\ref{imp:sec:update} explains why this is the relevant measure of quality).
The degraded provider is described in Section~\ref{imp:sec:degraded}. The student's row
gives its greedy accuracy on the 128-prompt selection subsample used during training.}
\label{imp:tab:pool}
\end{table}

The pool was designed so that the hypotheses can be tested rather than assumed. It
contains a strong and expensive provider (Phi-4), a mid-priced generalist
(Qwen2.5-7B), a code specialist at the same price (Qwen2.5-Coder-7B), a cheap provider
whose ceiling the student can plausibly reach during training (Qwen2.5-3B), and a cheap
provider from outside the student's model family (Phi-4-mini). The last of these
separates two properties that would otherwise coincide: without it, the only non-Qwen
provider would also be the only expensive one, and any effect of sharing a model family
with the student could not be told apart from an effect of price.

\subsection{Prices}
\label{imp:sec:prices}

Prices are set from measured compute and fixed before any training run. The price of
provider $k$ is the GPU time it needs to produce one response on one NVIDIA L40,
relative to Qwen2.5-7B:
\[
  c_k = \frac{\bar{\ell}_k / \nu_k}{\bar{\ell}_{\mathrm{ref}} / \nu_{\mathrm{ref}}},
\]
where $\bar{\ell}_k$ is the provider's mean response length in tokens over the
screening prompts of the three domains and $\nu_k$ is its generation throughput in tokens per second, both measured
during screening with the same serving software. A provider that writes long responses
is therefore more expensive per query, as it would be under per-token pricing. This
choice matters: an earlier estimate based on throughput alone had given Granite-3.3-8B a
price of 1.1, whereas its long responses put its measured price at 2.49.

Fixing prices in advance allows the switch points of Section~\ref{imp:sec:headroom} to
be stated as predictions before any training run, rather than fitted afterwards.

\subsection{Response cache}
\label{imp:sec:cache}

Following Section~\ref{sec:Problem Statement}, provider responses are generated once and
then charged as if they were live queries. Each provider answered every pool and
selection prompt $m = 4$ times, at temperature 0.7 with nucleus sampling at 0.95 and at
most 2,048 new tokens, using vLLM \autocite{kwon_vllm_2023}. This amounts to 30,000
responses per provider and took 8.1 GPU-hours for the five models. For the pool, only
the canonical answers are stored with each response; whether those answers are correct
is never recorded, so it cannot leak into training.

\subsection{The degraded provider}
\label{imp:sec:degraded}

The degraded provider is a copy of Phi-4 at the same price whose final answers are
corrupted on a fraction $p = 0.3$ of prompts. It tests whether a selector notices that a
provider is worse when its price gives no hint. Corruption is applied per prompt: with
probability $p$, decided by a hash of the prompt identifier, all four responses to the
prompt are replaced by one corrupted response. That response is the parent's response
carrying its most common answer, with the final answer changed to a different value:
another expression over the same integers for Countdown, a shifted value for MATH, and a
wrapper that perturbs the function's return values for code. No ground truth is used to
build it.

The per-prompt design follows from a short calculation. If each of the $m = 4$ responses
were corrupted independently with probability $p$, and corrupted answers were distinct,
the most common answer would survive whenever at least two responses survived, and the
provider's ceiling would fall only to about
\[
  \mu' \approx \mu \left(1 - p^4 - 4p^3(1 - p)\right) = 0.92\,\mu \quad \text{at } p = 0.3,
\]
which makes the degraded provider hard to distinguish from its parent. Corrupting whole
prompts instead gives $\mu' = (1 - p)\,\mu$ in expectation. Table~\ref{imp:tab:degraded} shows
that the built provider matches this prediction once the realised corruption rate is
used. On the screening data, per-response corruption at $p = 0.3$ lowered Phi-4's ceilings
only from 0.66, 0.76 and 0.74 to 0.56, 0.69 and 0.66.

\begin{table}[htbp]
\centering
\small
\begin{tabular}{@{}lccc@{}}
\toprule
 & \textbf{Countdown} & \textbf{MATH} & \textbf{Code} \\
\midrule
Phi-4 ceiling $\mu$                       & 0.874 & 0.804 & 0.686 \\
Prompts corrupted (selection set) $\hat{p}$ & 0.284 & 0.342 & 0.284 \\
Predicted ceiling $(1 - \hat{p})\,\mu$     & 0.626 & 0.529 & 0.491 \\
Measured ceiling of the degraded provider  & 0.614 & 0.532 & 0.492 \\
\bottomrule
\end{tabular}
\caption{The degraded provider against the per-prompt prediction.}
\label{imp:tab:degraded}
\end{table}

\section{Learning from a Black-Box Provider}
\label{imp:sec:update}

This section instantiates the operator $\operatorname{Update}$ of
Section~\ref{sec:Problem Statement}. The student is trained by black-box on-policy
distillation: it learns from its own responses, and the provider's responses are used
only to score them.

\subsection{Agreement reward}

For a prompt $x$ and provider $k$, let $y_{k,1}, \dots, y_{k,m}$ be the cached responses
and let
\[
  n_k(z; x) = \left|\{\, j : \kappa_d(y_{k,j}) = z \,\}\right|
\]
count how many of them give the canonical answer $z$. A student response $y$ with
answer $z = \kappa_d(y)$ receives the reward
\begin{equation}
  r_k(y; x) =
  \begin{cases}
    \dfrac{n_k(z; x)}{\max_{z'} n_k(z'; x)} & \text{if } z \ne \varnothing
      \text{ and the provider gave at least one valid answer,} \\[2ex]
    0 & \text{otherwise.}
  \end{cases}
  \label{imp:eq:reward}
\end{equation}
Matching the provider's most common answer earns a reward of 1, matching a less common
answer earns partial credit, and any other response earns 0. If the provider's answers to a prompt were $\{18, 18, 18, 22\}$, the student's answers
18, 22 and 7 would receive 1, $1/3$ and 0. This normalised vote share was chosen over two
alternatives: the indicator of matching the most common answer, and the unnormalised vote
share $n_k(z; x)/m$. Section~\ref{imp:sec:reward-evidence} gives the evidence. No ground
truth is involved, since the reward depends only on the student's answer and the
provider's answers.

\subsection{Policy update}

Each round produces $G = 8$ responses for each of the $b = 8$ prompts. Rewards are turned
into group-relative advantages as in GRPO (Chapter~\ref{chap:Preliminaries}),
\[
  A_{i,g} = \frac{r_{i,g} - \operatorname{mean}_{h}\, r_{i,h}}{\operatorname{std}_{h}\, r_{i,h} + \epsilon},
\]
with $A_{i,g} = 0$ for every response in a group whose rewards are all equal. The student
then takes one gradient step on
\begin{equation}
  \mathcal{L}(\theta) = -\frac{1}{\sum_{i,g} |y_{i,g}|}
    \sum_{i=1}^{b} \sum_{g=1}^{G} A_{i,g} \sum_{j=1}^{|y_{i,g}|}
    \log \pi_\theta\big(y_{i,g,j} \mid x_i, y_{i,g,<j}\big),
  \label{imp:eq:loss}
\end{equation}
where the sum over tokens is normalised by the total number of response tokens in the
batch, as in DAPO \autocite{yu_dapo_2025}. Because each batch of responses is used for
exactly one step, the importance ratio of GRPO equals one and its clipping never
applies, and no KL penalty or reference model is used. The update is therefore a plain
policy gradient with group-normalised rewards.

\subsection{Pre-query filter}

If all $G$ responses to a prompt share one canonical answer, including the case where
none has a valid answer, then every agreement reward gives them the same value and the
prompt contributes no gradient. Such prompts are dropped before the provider is
consulted, so they are neither charged nor trained on. The filter uses no labels and no
provider output, and on the screening data it removed 4 to 7\% of prompts in the three
domains without losing any group that would have carried a learning signal.

\subsection{What the student converges to}
\label{imp:sec:ceiling}

Ignoring the step size, the student's expected objective on a prompt $x$ is
$J(\pi) = \sum_{z} \pi(z \mid x)\, r_k(z; x)$, where $\pi(z \mid x)$ is the probability
that the student's answer is $z$. This is maximised by placing all probability on the
answer with the highest reward, which under Equation~\ref{imp:eq:reward} is the
provider's most common answer $\hat{z}_k(x)$. A student that fully learned from provider
$k$ would therefore be correct on exactly those prompts where the provider's most common
answer is correct. This motivates the ceiling
\[
  \mu_{k,d} = \Pr_{x \sim \mathcal{X}_d}\big[\, \hat{z}_k(x) \text{ is correct} \,\big],
\]
the provider's majority-vote accuracy \autocite{wang_selfconsistency_2023}, estimated on
the selection set and reported in Table~\ref{imp:tab:pool}. The ceiling, rather than the
provider's single-response accuracy, is the measure of provider quality used throughout.
The two can differ substantially: on Countdown during screening, Phi-4 answered 53\% of
prompts correctly per response but 66\% by majority vote.

The quality of the learning signal can be stated in the same terms. Let $q$ be the
expected reward of a correct response and $\varepsilon_z$ that of a wrong response with
answer $z$. If $\varepsilon_z = \varepsilon$ for every wrong answer, then
\[
  J(\pi) = \varepsilon + (q - \varepsilon) \Pr_\pi[\text{correct}],
  \qquad \nabla J = (q - \varepsilon)\, \nabla \Pr_\pi[\text{correct}],
\]
so learning from a provider is learning from correctness, slowed by the factor
$q - \varepsilon$, called the gap below. When $\varepsilon_z$ varies, the remainder pulls
the student towards the particular wrong answers that the provider favours. Separately, a
group of $G$ responses produces a gradient only if its rewards are not all equal. For a
binary reward with success probability $p$ per response this happens with probability
\[
  1 - p^G - (1 - p)^G,
\]
which is 0.22 for $p = 0.03$ and $G = 8$, close to the 0.20 measured for the base
student on Countdown with three and four numbers. The useful signal in a batch is
therefore proportional to the product of the fraction of informative groups and the gap.

\subsection{Evidence for the choice of reward}
\label{imp:sec:reward-evidence}

Table~\ref{imp:tab:reward} compares the three agreement rewards with the ground-truth
correctness reward on the base student's own responses, using the product of the
informative fraction and the gap. On MATH and code, both vote-share rewards carry about
as much signal as ground truth or more. The student's wrong answers tend to disagree with
one another, so agreement with a provider varies more within a group than correctness
does. On Countdown, the unnormalised vote share loses much of its gap, because providers
there often give some invalid answers, which keeps the vote share of even a correct
answer low. The normalised form removes that loss while keeping the higher informative
fraction, and it is the best or close to the best agreement reward in every column.

\begin{table}[htbp]
\centering
\small
\begin{tabular}{@{}llccc@{}}
\toprule
\textbf{Reward} & \textbf{Provider} & \textbf{Countdown (3 and 4)} & \textbf{MATH} & \textbf{Code} \\
\midrule
Correctness (ground truth) & --         & 0.20 & 0.57 & 0.48 \\
\midrule
Most common answer         & Qwen2.5-7B & 0.14 & 0.59 & 0.42 \\
                           & Phi-4      & 0.11 & 0.57 & 0.42 \\
Vote share                 & Qwen2.5-7B & 0.12 & 0.65 & 0.50 \\
                           & Phi-4      & 0.12 & 0.61 & 0.47 \\
Normalised vote share      & Qwen2.5-7B & 0.18 & 0.64 & 0.53 \\
                           & Phi-4      & 0.13 & 0.61 & 0.48 \\
\bottomrule
\end{tabular}
\caption{Learning signal per batch (fraction of informative groups multiplied by the
gap) for each reward, on 200 screening prompts per domain with $G = 8$ responses from the
base student. Countdown here includes four-number puzzles.}
\label{imp:tab:reward}
\end{table}

The cost of partial credit is a slightly higher reward for wrong answers: for
Qwen2.5-7B on MATH and code, $\varepsilon$ is 0.08 under the normalised vote share against
0.05 under the most-common-answer indicator. The value of $\varepsilon$ is logged during
training (Section~\ref{imp:sec:groundtruth}).

\section{Selection Policies and Baselines}
\label{imp:sec:selectors}

\subsection{Blocks and feedback}

Selectors choose an arm once per block of $L = 10$ rounds rather than every round. One
gradient step changes selection accuracy by well under one percentage point, whereas an
evaluation on 128 prompts has a standard error of about four percentage points, so
feedback after a single round would be almost pure noise. At the end of each block the
student is evaluated greedily, at temperature zero, on a fixed subsample of 128
selection prompts per domain. This yields the context
$s = (\alpha_1, \dots, \alpha_D)$ of Section~\ref{sec:Problem Statement}, the student's
current accuracy in each domain. A selector only considers arms it can still afford for
at least one round.

\subsection{Headroom-per-cost selector}
\label{imp:sec:headroom}

The proposed selector estimates the marginal value $v(k, d; \theta)$ of
Section~\ref{sec:Problem Statement} with a simple model. Section~\ref{imp:sec:ceiling}
showed that learning from provider $k$ drives the student towards the ceiling
$\mu_{k,d}$ rather than beyond it. If progress per round is further assumed to be proportional to the remaining
headroom, the value of an arm per unit of cost is
\[
  \hat{v}(k, d) = \frac{\max\big(\mu_{k,d} - \alpha_d,\ 0\big)}{c_k}.
\]
The selector plays the arm of highest $\hat{v}$, and with probability 0.1 a uniformly
random affordable arm instead, so that every arm continues to be observed. The ceilings
are estimated once from the providers' cached responses to the selection set, which is
labelled and may be used for this purpose. A label-free alternative, scoring each
provider by how often it agrees with the other providers, was examined and rejected
because it favoured the model family that made up most of the pool
(Chapter~\ref{chap:Evaluation}); this is the conditional-independence assumption of
\textcite{dawid_skene_1979} failing in practice.

The model makes the hypotheses concrete. A provider $e$ that is more expensive than a
provider $c$ but has a higher ceiling becomes the better choice once the student's
accuracy exceeds
\[
  \alpha^{*} = \frac{c_e\, \mu_{c,d} - c_c\, \mu_{e,d}}{c_e - c_c},
\]
the point at which $(\mu_{e,d} - \alpha)/c_e$ overtakes $(\mu_{c,d} - \alpha)/c_c$.
Table~\ref{imp:tab:predictions} lists what the model predicts for the pool in
Table~\ref{imp:tab:pool}. The best arm at the start differs between domains, which is a
prediction of H1, and each domain has a point at which the best provider changes, which
is a prediction of H2. The prices and ceilings behind these predictions were fixed
before any training run, and Chapter~\ref{chap:Evaluation} tests the predictions.

\begin{table}[htbp]
\centering
\small
\begin{tabular}{@{}lccl@{}}
\toprule
\textbf{Domain} & \textbf{Student at start} & \textbf{Best arm at start} & \textbf{Predicted first switch} \\
\midrule
Countdown & 0.195 & Qwen2.5-3B       & to Qwen2.5-7B at $\alpha^{*} = 0.27$ \\
MATH      & 0.414 & Qwen2.5-3B       & to Qwen2.5-7B at $\alpha^{*} = 0.63$ \\
Code      & 0.438 & Qwen2.5-Coder-7B & to Phi-4 at $\alpha^{*} = 0.57$ \\
\bottomrule
\end{tabular}
\caption{Predictions of the headroom model for the pool in Table~\ref{imp:tab:pool}.}
\label{imp:tab:predictions}
\end{table}

The assumption that progress is proportional to headroom is a modelling choice, not a
result. It ignores transfer between domains, differences in how easily the student
absorbs each provider's style, and the student's own capacity. The branching
measurements of Section~\ref{imp:sec:marginal} provide an independent check of it.

\subsection{Discounted Thompson sampling}

The model-free comparison treats each arm as a bandit arm whose payoff is the learning
progress it produces per unit of cost, as in automated curriculum learning
\autocite{graves_curriculum_2017,matiisen_tscl_2017}. After a block on arm $(k, d)$, the
selector observes
\[
  \rho = \frac{1}{D \, C_{\mathrm{blk}}} \sum_{d'=1}^{D}
         \big(\alpha^{\mathrm{after}}_{d'} - \alpha^{\mathrm{before}}_{d'}\big),
\]
where $C_{\mathrm{blk}}$ is what the block cost. This is the average accuracy gain over all
domains per unit of budget, so that transfer to other domains counts. It keeps a Gaussian posterior for each arm's mean payoff and plays the
arm with the highest sample, in the spirit of Thompson sampling for budgeted bandits
\autocite{xia_budgeted_ts_2015}. Because H2 states that payoffs drift as the student
learns, all past observations are discounted by a factor $\gamma = 0.9$ per block, which
is the approach of discounted bandits for changing environments
\autocite{garivier_switching_2008,raj_dts_2017}. The posterior's noise scale is set to
the spread of the arms' current mean payoffs.

\subsection{Baselines}
\label{imp:sec:baselines}

\begin{description}[itemsep=3pt,topsep=4pt]
  \item[Uniform.] Cycles through all $K \times D$ arms in a fixed random order, giving
    each arm the same number of blocks.
  \item[Single provider.] Uses one provider for the whole run, cycling through the
    domains. Run once for each of the five real providers, it covers the
    strongest-provider-only baseline (Phi-4) and the cheapest-provider-only baseline
    (Qwen2.5-3B). The best of the five on the selection set serves as the best static
    allocation in hindsight. This restricts static allocations to a single provider,
    which Section~\ref{imp:sec:limitations} discusses.
  \item[Correctness reward without providers.] The same training loop with the
    agreement reward replaced by ground-truth correctness on the pool, no provider and no
    cost, run for 200 rounds with domains in rotation. It shows what the student could
    learn if labels were available, without any teacher.
  \item[Supervised fine-tuning.] An off-policy baseline. It spends the same budget $B$
    on provider responses, divided equally over all arms. For each bought prompt it
    keeps the response that carries the provider's most common answer, a label-free filter
    in the spirit of self-consistency \autocite{wang_selfconsistency_2023}. It then
    fine-tunes the same LoRA configuration on these responses for two epochs.
\end{description}

The greedy oracle listed among the objectives in Section~\ref{sec:Objectives} is
approximated with the measurements of the next section: at each checkpoint it would
choose the arm whose measured marginal value is highest. It serves as a reference for
the selectors' choices rather than as a separate training run.

\section{Measuring Marginal Value}
\label{imp:sec:marginal}

H1 and H2 concern the marginal value $v(k, d; \theta)$, which is measured directly by
branching. The uniform-allocation run of each seed provides checkpoints at 25, 50 and
75\% of its budget, and the untrained student serves as the checkpoint at 0\%. From each checkpoint $\theta_c$, a separate short run trains on a
single arm $(k, d)$ for $R = 20$ rounds and is then evaluated on the full selection set.
The estimate of marginal value is
\[
  \hat{v}(k, d; \theta_c) =
  \frac{U(\theta_c^{k,d}) - U(\theta_c)}{\text{budget spent in the branch}},
\]
where $\theta_c^{k,d}$ is the student after the branch and $U$ is the mean selection
accuracy over the three domains. The per-domain change is also recorded. The uniform run
is used as the source of checkpoints because it visits every arm and therefore favours
none.

Branches from the same checkpoint and seed share everything except the arm: they start
from the same adapter and optimiser state, read the same prompt sequence for a given
domain, and draw student responses with the same random seeds. Differences between two
providers on the same domain are therefore paired comparisons, which reduces their
variance. With $K = 6$ providers including the degraded one, $D = 3$ domains, $C = 4$
checkpoints and $S = 3$ seeds, the design has 216 branches.

H1 is supported if, at a fixed checkpoint, the provider with the highest $\hat{v}$
differs between domains. H2 is supported if, within a domain, the provider with the
highest $\hat{v}$ changes between checkpoints. H3 is tested by comparing the final test
utility of the adaptive selectors with that of the best single-provider run at equal
budget.

\section{Experimental Configuration}
\label{imp:sec:config}

Table~\ref{imp:tab:config} lists the settings shared by all runs. Values marked with an
asterisk were set before the pilot run of Chapter~\ref{chap:Evaluation} and may be
revised in light of it.

\begin{table}[htbp]
\centering
\small
\begin{tabular}{@{}lll@{}}
\toprule
\textbf{Component} & \textbf{Setting} & \textbf{Value} \\
\midrule
Budget         & total budget $B$                         & 2,400 price units \\
               & prompts per round $b$                    & 8 \\
               & round cap $T_{\max}$                     & 400 \\
Rollouts       & responses per prompt $G$                 & 8 \\
               & temperature, top-$p$                     & 1.0, 1.0 \\
               & maximum response length                  & 1,024 tokens \\
Optimisation   & optimiser                                & AdamW, no weight decay \\
               & learning rate$^{*}$                      & $2 \times 10^{-5}$, constant \\
               & gradient clipping                        & 1.0 \\
               & LoRA rank, scaling $\alpha$, dropout     & 16, 32, 0 \\
Selection      & block length $L$                         & 10 rounds \\
               & selection subsample per domain           & 128 prompts, greedy decoding \\
               & exploration (headroom selector)          & 0.1 \\
               & discount $\gamma$ (Thompson sampling)    & 0.9 \\
Providers      & cached responses per prompt $m$          & 4 \\
               & sampling for the cache                   & temperature 0.7, top-$p$ 0.95, 2,048 tokens \\
Replication    & seeds per condition                      & 3 \\
\bottomrule
\end{tabular}
\caption{Configuration shared by all runs.}
\label{imp:tab:config}
\end{table}

\paragraph{Budget.} The budget is set so that the uniform policy, whose average price per
prompt is the mean price of 1.53, runs for about 200 rounds. At equal budget, cheaper
providers buy more training steps: a run using only Phi-4 stops after about 103 rounds,
while a run using only Qwen2.5-3B would be allowed 600 and is stopped by the cap at 400.
Results are therefore reported both against budget spent and against the number of
gradient steps, so that a selector's gains can be separated from the extra steps that
cheap providers buy. Checkpoints are saved when a run's progress, the larger of its
budget fraction and its round fraction, first reaches 25, 50 and 75\%.

\paragraph{Prompt order.} Each domain's pool is read in a fixed random order per seed, so
no prompt is repeated until the pool is exhausted. Student responses are drawn with a
seed derived from the run seed and the round number, which gives common random numbers
across runs with the same seed.

\paragraph{Ground truth.} \label{imp:sec:groundtruth}Labels enter a run in exactly three
places: the evaluations on the selection set, the correctness reward of the
provider-free baseline, and diagnostic quantities written to the log, such as the
student's accuracy on the pool and the reward given to correct and wrong responses
($q$ and $\varepsilon$). No decision made during a run reads the diagnostics.

\paragraph{Hardware and software.} Experiments run on a shared server with six NVIDIA
L40 GPUs of 48\,GB each, without a job scheduler. Jobs wait for an idle GPU and use at
most two at a time. Each run occupies one GPU, because two runs sharing a GPU were
measured to be slower in total than the same runs in sequence. Within a run, vLLM
\autocite{kwon_vllm_2023} generates the student's responses with the current LoRA adapter,
which is reloaded after every update, and PyTorch with the PEFT library computes the
update. Scoring runs in parallel worker processes, one per response. In the pilot, a
round took a median of about 10 seconds on Countdown and 14 seconds on MATH.

\section{Verification of the Testbed}
\label{imp:sec:verification}

Each component was checked before the main runs.

\begin{description}[itemsep=3pt,topsep=4pt]
  \item[Answer checkers.] Math-Verify accepted all 300 reference solutions in a sample of
    MATH problems. In the code domain, only problems whose reference solution passes in
    the sandbox are used (92.5\% of sampled candidates).
  \item[Sandbox.] Test programs confirmed that writes outside the working directory fail,
    network access fails, and infinite loops are stopped by the time limits.
  \item[Harness rules.] Each rule of Section~\ref{imp:sec:sandbox} was introduced in
    response to a measured failure and applies to every model. Removing appended
    self-tests recovered 286 of the base student's screening responses whose own failing
    test had caused them to be scored as missing answers. Normalising LaTeX in Countdown
    answers raised Phi-4's screening accuracy from 0.35 to 0.52 and changed no other
    model's by more than 0.01. Exposing the solution's names to test files that
    omit their import raised the pass rate of KodCode reference solutions from about 80\%
    to 92.5\%.
  \item[Degraded provider.] Its measured ceilings match the per-prompt prediction
    (Table~\ref{imp:tab:degraded}).
  \item[Training loop.] Short end-to-end runs confirmed budget accounting, the pre-query
    filter, checkpointing, resumption of a branch from a checkpoint, and the provider-free
    baseline. In the selector tests, the headroom selector moved from Qwen2.5-3B on MATH to
    Phi-4 on code as the simulated student accuracy rose, and Thompson sampling learned to
    prefer the arm with the highest simulated payoff.
\end{description}

\section{Limitations of the Design}
\label{imp:sec:limitations}

\begin{description}[itemsep=3pt,topsep=4pt]
  \item[Cached responses.] Each provider has four responses per prompt, generated once.
    The student is thus trained against a fixed sample of each provider's behaviour, and
    live querying would add variation that is not modelled.
  \item[Prices.] Prices are a compute-based proxy measured on one type of GPU with one
    serving system. Real API prices differ, and the switch points of
    Table~\ref{imp:tab:predictions} depend on price ratios.
  \item[Model families.] Three of the five real providers and the student are from the
    Qwen family. Phi-4-mini was added to separate family from price, but the pool remains
    unbalanced.
  \item[Agreement in the code domain.] Canonical answers for code are the function's
    outputs on the test inputs. For KodCode, outputs are recorded up to the first failing
    assertion in each test function, so the expected outputs influence which calls are
    recorded, although not the values compared.
  \item[Degraded provider.] Corrupted code answers change the function's return values.
    For about 5\% of corrupted code prompts this does not change the tested behaviour, so
    the corruption is slightly weaker in that domain.
  \item[Static allocations.] The best static allocation in hindsight is taken over
    single-provider runs only. A static mixture of providers chosen per domain could do
    better and is not searched.
  \item[The headroom model.] Its assumption that progress is proportional to headroom is
    a simplification, and the branching measurements are the test of whether it holds.
  \item[Scale.] There is one student of 1.5B parameters, three domains and three seeds
    per condition, chosen to fit the available compute.
\end{description}
